% !TeX root = ../../main.tex
\section{最小作用原理和拉格朗日力学}
\begin{solution}[谐振子的拉格朗日量]
\problem{考虑以下被称为谐振子的系统。质量为 \(m\) 的物块位于一个光滑表面上，弹簧的劲度系数为 \(k\)。写出拉格朗日量并导出运动方程。}
\begin{center}
    \includegraphics[width=0.55\textwidth]{harmonic_oscillator.png}
\end{center}


\answer
取物块沿弹簧伸缩方向的位移为广义坐标 \(x\)，并约定 \(x=0\) 为弹簧处于自然长度时物块的位置。由于物块只在水平方向运动，且表面光滑无摩擦，所以系统只有一个自由度，广义速度为 \(\dot{x}\)。物块的动能为
\[
T=\frac{1}{2}m\dot{x}^{2}.
\]
弹簧的弹性势能只与弹簧相对自然长度的伸长量或压缩量有关，因此为
\[
V=\frac{1}{2}kx^{2}.
\]
由拉格朗日量的定义 \(\mathcal{L}=T-V\)，可得该谐振子的拉格朗日量为
\[
\mathcal{L}(x,\dot{x})=\frac{1}{2}m\dot{x}^{2}-\frac{1}{2}kx^{2}.
\]
下面由欧拉--拉格朗日方程导出运动方程。对广义速度求偏导，有
\[
\frac{\partial \mathcal{L}}{\partial \dot{x}}=m\dot{x},
\]
再对时间求导，得到
\[
\frac{d}{dt}\left(\frac{\partial \mathcal{L}}{\partial \dot{x}}\right)=m\ddot{x}.
\]
另一方面，对广义坐标求偏导，有
\[
\frac{\partial \mathcal{L}}{\partial x}=-kx.
\]
将以上结果代入欧拉--拉格朗日方程
\[
\frac{d}{dt}\left(\frac{\partial \mathcal{L}}{\partial \dot{x}}\right)-\frac{\partial \mathcal{L}}{\partial x}=0,
\]
得到
\[
m\ddot{x}-(-kx)=0,
\]
即
\[
m\ddot{x}+kx=0.
\]
因此该系统的运动方程为
\[
m\ddot{x}=-kx,
\]
或者等价地写为
\[
\ddot{x}=-\frac{k}{m}x.
\]
这表明物块所受的回复力与位移大小成正比、方向与位移方向相反，系统作角频率为 \(\omega=\sqrt{k/m}\) 的简谐振动。
\end{solution}


\begin{solution}[耦合质量问题的拉格朗日量]
\problem{对于第 1.8 节末尾讨论的耦合质量问题做同样处理。把得到的运动方程与 \((1.8.24)\) 和 \((1.8.25)\) 式进行比较。}
\answer
设两个质量相同的物块质量均为 \(m\)，它们只能沿同一直线运动。取两个物块相对于各自平衡位置的位移分别为广义坐标 \(x_1\) 和 \(x_2\)，并取向右为正方向。由于系统中有两个独立的广义坐标，因此拉格朗日量应写成 \(\mathcal{L}(x_1,x_2,\dot{x}_1,\dot{x}_2)\)。两个物块的动能之和为
\[
T=\frac{1}{2}m\dot{x}_1^2+\frac{1}{2}m\dot{x}_2^2=\frac{1}{2}m(\dot{x}_1^2+\dot{x}_2^2).
\]
势能由三根弹簧贡献：左侧弹簧的形变量为 \(x_1\)，右侧弹簧的形变量为 \(x_2\)，中间弹簧的相对形变量为 \(x_2-x_1\)。因此系统的总势能为
\[
V=\frac{1}{2}k\left[x_1^2+x_2^2+(x_2-x_1)^2\right].
\]
由 \(\mathcal{L}=T-V\)，得到该耦合质量系统的拉格朗日量为
\[
\mathcal{L}=\frac{1}{2}m(\dot{x}_1^2+\dot{x}_2^2)-\frac{1}{2}k\left[x_1^2+x_2^2+(x_2-x_1)^2\right].
\]
下面分别对 \(x_1\) 和 \(x_2\) 使用欧拉--拉格朗日方程
\[
\frac{d}{dt}\left(\frac{\partial \mathcal{L}}{\partial \dot{x}_i}\right)-\frac{\partial \mathcal{L}}{\partial x_i}=0,\qquad i=1,2.
\]
先对 \(x_1\) 计算。由于
\[
\frac{\partial \mathcal{L}}{\partial \dot{x}_1}=m\dot{x}_1,\qquad \frac{d}{dt}\left(\frac{\partial \mathcal{L}}{\partial \dot{x}_1}\right)=m\ddot{x}_1,
\]
并且
\[
\frac{\partial \mathcal{L}}{\partial x_1}
=-\frac{1}{2}k\frac{\partial}{\partial x_1}\left[x_1^2+x_2^2+(x_2-x_1)^2\right]
=-\frac{1}{2}k\left[2x_1-2(x_2-x_1)\right]
=-k(2x_1-x_2),
\]
所以
\[
m\ddot{x}_1-\big[-k(2x_1-x_2)\big]=0,
\]
即
\[
m\ddot{x}_1=-2kx_1+kx_2.
\]
再对 \(x_2\) 计算。由于
\[
\frac{\partial \mathcal{L}}{\partial \dot{x}_2}=m\dot{x}_2,\qquad \frac{d}{dt}\left(\frac{\partial \mathcal{L}}{\partial \dot{x}_2}\right)=m\ddot{x}_2,
\]
并且
\[
\frac{\partial \mathcal{L}}{\partial x_2}
=-\frac{1}{2}k\frac{\partial}{\partial x_2}\left[x_1^2+x_2^2+(x_2-x_1)^2\right]
=-\frac{1}{2}k\left[2x_2+2(x_2-x_1)\right]
=-k(2x_2-x_1),
\]
所以
\[
m\ddot{x}_2-\big[-k(2x_2-x_1)\big]=0,
\]
即
\[
m\ddot{x}_2=kx_1-2kx_2.
\]
因此，由拉格朗日方程得到的运动方程为
\[
\begin{cases}
m\ddot{x}_1=-2kx_1+kx_2,\\
m\ddot{x}_2=kx_1-2kx_2.
\end{cases}
\]
两式同除以 \(m\)，可写为
\[
\begin{cases}
\ddot{x}_1=-\dfrac{2k}{m}x_1+\dfrac{k}{m}x_2,\\[6pt]
\ddot{x}_2=\dfrac{k}{m}x_1-\dfrac{2k}{m}x_2.
\end{cases}
\]
这正是由牛顿第二定律得到的耦合质量系统运动方程，与 \((1.8.24)\) 和 \((1.8.25)\) 式一致。其一致性说明，在取相同广义坐标和相同势能零点的条件下，拉格朗日方法与牛顿力学方法给出的动力学描述完全等价。
\end{solution}


\begin{solution}[拉格朗日量，极坐标]
\problem{一个在三维空间中运动的质量为 \(m\) 的粒子受到势 \(V(r,\theta,\phi)=V(r)\) 的作用。写出它的 \(\mathcal{L}\)，并求其运动方程。}
\answer
由于势能只依赖于径向坐标 \(r\)，这是一个中心势问题。采用球坐标 \((r,\theta,\phi)\)，其中
\[
x=r\cos\theta\sin\phi,\qquad y=r\sin\theta\sin\phi,\qquad z=r\cos\phi .
\]
在这种约定下，速度平方为
\[
\dot{x}^{2}+\dot{y}^{2}+\dot{z}^{2}
=\dot{r}^{2}+r^{2}\dot{\phi}^{2}+r^{2}\dot{\theta}^{2}\sin^{2}\phi .
\]
因此粒子的动能为
\[
T=\frac{1}{2}m\left(\dot{r}^{2}+r^{2}\dot{\phi}^{2}+r^{2}\dot{\theta}^{2}\sin^{2}\phi\right),
\]
势能为 \(V=V(r)\)，所以拉格朗日量为
\[
\mathcal{L}
=\frac{1}{2}m\left(\dot{r}^{2}+r^{2}\dot{\phi}^{2}+r^{2}\dot{\theta}^{2}\sin^{2}\phi\right)-V(r).
\]
下面分别对广义坐标 \(r,\phi,\theta\) 使用欧拉--拉格朗日方程
\[
\frac{d}{dt}\left(\frac{\partial \mathcal{L}}{\partial \dot{q}}\right)-\frac{\partial \mathcal{L}}{\partial q}=0.
\]
先对 \(r\) 计算。由于
\[
\frac{\partial \mathcal{L}}{\partial \dot{r}}=m\dot{r},\qquad
\frac{d}{dt}\left(\frac{\partial \mathcal{L}}{\partial \dot{r}}\right)=m\ddot{r},
\]
并且
\[
\frac{\partial \mathcal{L}}{\partial r}
=mr\dot{\phi}^{2}+mr\dot{\theta}^{2}\sin^{2}\phi-\frac{dV}{dr},
\]
所以径向运动方程为
\[
m\ddot{r}=mr\dot{\phi}^{2}+mr\dot{\theta}^{2}\sin^{2}\phi-\frac{dV}{dr}.
\]
再对 \(\phi\) 计算。由于
\[
\frac{\partial \mathcal{L}}{\partial \dot{\phi}}=mr^{2}\dot{\phi},\qquad
\frac{d}{dt}\left(\frac{\partial \mathcal{L}}{\partial \dot{\phi}}\right)
=m\left(r^{2}\ddot{\phi}+2r\dot{r}\dot{\phi}\right),
\]
并且
\[
\frac{\partial \mathcal{L}}{\partial \phi}
=mr^{2}\dot{\theta}^{2}\sin\phi\cos\phi,
\]
所以 \(\phi\) 方向的运动方程为
\[
m\left(r^{2}\ddot{\phi}+2r\dot{r}\dot{\phi}\right)
=mr^{2}\dot{\theta}^{2}\sin\phi\cos\phi.
\]
最后对 \(\theta\) 计算。拉格朗日量中不显含 \(\theta\)，因此 \(\theta\) 是循环坐标，有
\[
\frac{\partial \mathcal{L}}{\partial \theta}=0.
\]
又因为
\[
\frac{\partial \mathcal{L}}{\partial \dot{\theta}}
=mr^{2}\dot{\theta}\sin^{2}\phi,
\]
所以
\[
\frac{d}{dt}\left(mr^{2}\dot{\theta}\sin^{2}\phi\right)=0.
\]
综上，粒子在中心势 \(V(r)\) 中的运动方程为
\[
\begin{cases}
m\ddot{r}=mr\dot{\phi}^{2}+mr\dot{\theta}^{2}\sin^{2}\phi-\dfrac{dV}{dr},\\[6pt]
m\left(r^{2}\ddot{\phi}+2r\dot{r}\dot{\phi}\right)=mr^{2}\dot{\theta}^{2}\sin\phi\cos\phi,\\[6pt]
\dfrac{d}{dt}\left(mr^{2}\dot{\theta}\sin^{2}\phi\right)=0.
\end{cases}
\]
其中第三个方程表明，与循环坐标 \(\theta\) 对应的广义动量
\(
p_{\theta}=mr^{2}\dot{\theta}\sin^{2}\phi
\)
守恒。这正反映了中心势问题的旋转对称性。
\end{solution}